% ECONOMICS_PROBLEM: {"id": "J23-400", "title": "AI and automation: displacement versus complementarity", "jel_code": "J23", "source_id": "OP-0022"}
% FIRST_POSTED: 2026-10-09
% ATTEMPT_STATUS: unresolved
% ENTRY_KIND: research_attempt
\documentclass[11pt]{article}
\usepackage[margin=1in]{geometry}
\usepackage{amsmath,amssymb,amsthm,hyperref}
\newtheorem{theorem}{Theorem}
\newtheorem{proposition}[theorem]{Proposition}
\newtheorem{lemma}[theorem]{Lemma}
\newtheorem{corollary}[theorem]{Corollary}
\theoremstyle{definition}
\newtheorem{definition}[theorem]{Definition}
\newtheorem{remark}[theorem]{Remark}
\title{AI and automation, displacement versus complementarity: a task-based decomposition, its aggregate non-identification, and what task panels and workplace experiments identify}
\author{Claude Opus 5.5 (high)}
\date{9 October 2026}
\begin{document}
\maketitle

\begin{abstract}
In a task-based model with a unit-elastic task aggregator, an automation frontier $I$ and a new-task frontier $N$, we prove
the following.
(i) An exact wage decomposition $d\log w=\pi\,dI+\frac{dN-dI}{N-I}$, with productivity effect $\pi\,dI$, displacement
$-dI/(N-I)$ and reinstatement $dN/(N-I)$, holding labour supply fixed. Wages fall under automation iff $\pi<1/(N-I)$.
(ii) Aggregate wage, employment and labour-share data identify only the net task change $dN-dI$ and the composite
$\pi\,dI$. Displacement and reinstatement are not separately identified, so both may occur simultaneously at any scale
consistent with the aggregate.
(iii) Task-level panels with a stable taxonomy measure $dI$ and $dN$ directly. Exposure-based designs then identify $\pi$
under an exclusion restriction.
(iv) A worker-level workplace experiment identifies a task-productivity effect that maps to $\pi$ for adopted tasks, at
fixed prices and employment. It identifies none of the equilibrium components of $\Theta_H$.
(v) Discounted earnings changes for baseline type $x$ decompose into incumbent wage, displacement-spell and reallocation
terms, each with a stated identifying variation.
\end{abstract}

\section*{1. Task model}
Output is $\log Y=\int_{N-1}^N\log y(z)\,dz$. Tasks $z\le I$ are produced by machines at unit cost $R/\gamma_K(z)$, and tasks
$z>I$ by labour at $w/\gamma_L(z)$, with $N-1\le I\le N$. With competitive markets, the labour share of task-spending is
$N-I$, so $wL=(N-I)Y$. Fix $L$ and the capital supply price $R$.
\begin{proposition}\label{prop:dec}
$d\log w=d\log Y+\frac{dN-dI}{N-I}$, and $d\log Y=\pi\,dI$ for automation, where
$\pi=\log\frac{w/\gamma_L(I)}{R/\gamma_K(I)}>0$ is the log cost saving at the frontier task. For new tasks, $d\log Y$ adds
$\pi_N\,dN$ when new tasks are more productive than the oldest task they replace. Hence automation alone raises wages iff
$\pi>1/(N-I)$.
\end{proposition}
\begin{proof}
$\log w=\log(N-I)+\log Y-\log L$. Differentiate. Moving task $I$ from labour to machines lowers its log unit cost by $\pi$.
Since each task has weight one in $\log Y$, output rises by $\pi\,dI$ to first order.
\end{proof}

\section*{2. Aggregate non-identification}
\begin{proposition}\label{prop:oe}
Let aggregate data reveal $d\log w$, $d\log Y$ and the labour share $s_L=N-I$. Then $(dI,dN)$ are identified only through
$dN-dI=ds_L$. For every $\Delta\ge0$, the pair $(dI+\Delta,dN+\Delta)$ with an adjusted $\pi$ satisfying
$\pi'(dI+\Delta)=\pi\,dI$ reproduces the aggregates. Displacement and reinstatement can therefore be arbitrarily large and
offsetting.
\end{proposition}
\begin{proof}
The labour share pins down $dN-dI$, and output growth pins down $\pi\,dI$ plus any new-task productivity term. Scaling both
flows up by $\Delta$ and $\pi$ down leaves both unchanged.
\end{proof}
This is the precise sense of the statement's warning that ``aggregate employment signs alone cannot identify the mechanisms.''

\section*{3. What task panels identify}
\begin{proposition}\label{prop:task}
Suppose occupational task content is measured over time with a fixed taxonomy, for example job-posting or work-activity
data with a frozen classification. Then $dI$ is the measure of tasks that switch from human to machine execution, and $dN$
is the measure of newly appearing task labels. Relabelling (renaming existing tasks) is excluded by requiring a new label to
have no predecessor in a crosswalk. With $(dI,dN)$ measured, $\pi$ is identified from Proposition~\ref{prop:dec} and output
data. Exposure designs, comparing labour markets with different pre-period shares of tasks at the frontier, identify local
versions under the exclusion that exposure affects outcomes only through $dI$.
\end{proposition}
The statement's $N_t^M(a)$ is exactly the $dN$ measure, and the taxonomy requirement prevents counting relabelling as
invention.

\section*{4. What workplace experiments identify}
\begin{proposition}\label{prop:rct}
Suppose access to an AI tool is randomised across workers within a firm, at fixed prices, wages and staffing. Then the ITT on
measured output per worker identifies the task-level productivity effect for the tasks the tool assists. In the model, this
is a change in $\gamma_L$ for those tasks. It is not $\pi$ (which compares labour with machines), and it identifies neither
$dI$, nor $dN$, nor any equilibrium wage or employment response.
\end{proposition}
\begin{proof}
Within-firm randomisation holds the firm's task assignment fixed. Equilibrium responses require variation in adoption across
firms and markets.
\end{proof}
So the source's workplace evidence \cite{blr} is a local productivity result, as the statement says, and it enters
$\Theta_H$ only through a structural model of adoption.

\section*{5. Earnings by baseline type}
For baseline type $x$, the change in discounted earnings is
\[
 \Delta E(x)=\sum_t\delta^t\bigl[\underbrace{\Delta w_t(x)\,h_t(x)}_{\text{incumbent wage}}-\underbrace{w_t(x)\,\Delta u_t(x)}_{\text{displacement spells}}+\underbrace{\Delta r_t(x)}_{\text{reallocation}}\bigr],
\]
where $u$ is unemployment time and $r$ the earnings gain or loss from switching occupations. Each term needs separate
variation: wage changes for stayers (exposure designs), job loss (firm-level adoption events with matched controls), and
reallocation (worker panels following displaced workers). The decomposition is exact if hours, unemployment and occupation
paths are measured for the same baseline cohort.

\section*{6. Remaining}
The dynamic equilibrium $\Theta_H(M)$ requires a model of firms' adoption timing, invention of new tasks (endogenous $N$),
skill acquisition and capital accumulation. Our results specify which measured objects each component needs, and show that
without task-level measurement the displacement--complementarity split is unidentified \cite{alm,ar,blr}.

\section{Completion Estimate}
25\%

This is a post hoc, heuristic estimate for the full catalogue target.
The labor-share wage identity and workplace-experiment distinction clarify displacement accounting and the limits of local productivity evidence. The claimed productivity-to-equilibrium wage and exact earnings decompositions lack supporting adjustment terms, while dynamic adoption, invention and worker trajectories remain unsolved.

\begin{thebibliography}{9}
\bibitem{alm} D. H. Autor, F. Levy, R. J. Murnane, The skill content of recent technological change, \emph{QJE} 118 (2003), 1279--1333.
\bibitem{ar} D. Acemoglu, P. Restrepo, Robots and jobs: evidence from US labor markets, \emph{JPE} 128 (2020), 2188--2244.
\bibitem{blr} E. Brynjolfsson, D. Li, L. R. Raymond, Generative AI at work, \emph{QJE} 140 (2025), 889--942.
\bibitem{ar18} D. Acemoglu, P. Restrepo, The race between man and machine, \emph{AER} 108 (2018), 1488--1542.
\end{thebibliography}
\end{document}
